\documentclass[11pt]{article}
\usepackage[margin=1in]{geometry}
\usepackage{amsmath,amssymb,amsthm}
\usepackage{graphicx}
\usepackage{hyperref}

\newtheorem{theorem}{Theorem}
\newtheorem{lemma}{Lemma}
\newtheorem{corollary}{Corollary}
\newtheorem{remark}{Remark}

\newcommand{\F}{\mathcal F}
\newcommand{\M}{\mathcal M}
\newcommand{\Eext}{E_{\mathrm{ext}}}
\newcommand{\diam}{\operatorname{diam}}

\title{A One-Point Rigidification for Finite Lipschitz-Free Spaces}
\author{Run \texttt{fa\_banach\_001}, lane 18}
\date{2026-07-03}

\begin{document}
\maketitle

\section*{Source And Open Problem}

The source paper is Cuth--Doucha--Titkos, \emph{Isometries of Lipschitz-free Banach spaces}, arXiv:2402.08266 \cite{CDT}.  In Section 8.2 it asks:

\begin{center}
\includegraphics[width=0.96\textwidth]{figures/open_problem_crop.png}
\end{center}

We address Question 3 in a finite graph-theoretic subcase.  The source proves that every metric space embeds isometrically into a Lipschitz-free rigid metric space after adding three points.  It asks whether one point always suffices.

\section*{Proof Intuition}

The source's main rigidity criterion says that a weak Prague metric space is Lipschitz-free rigid once its preserved-extreme graph is 3-connected.  If a finite metric space already has a 2-connected preserved-extreme graph, then adding one far-away point joined at the same distance to every old point turns that graph into its cone.  Coning raises 2-connectivity to 3-connectivity.  The large common distance prevents the new point from lying between old pairs, while every new radial pair becomes preserved extreme.

\section*{Result}

For a finite metric space \(M\), let \(G_{\mathrm{ext}}(M)\) be the undirected graph whose vertices are the endpoints of preserved extreme molecules and whose edges are the unordered pairs \(\{x,y\}\) such that \(m_{x,y}\) is a preserved extreme point of the unit ball of \(\F(M)\).

\begin{theorem}[One-point cone rigidification]
\label{thm:one-point}
Let \(M\) be a finite metric space such that \(G_{\mathrm{ext}}(M)\) has vertex set \(M\) and is 2-connected.  Fix \(R>\diam(M)\), add a point \(p\notin M\), and define a metric \(d_R\) on
\[
  M_R=M\cup\{p\}
\]
by keeping the original metric on \(M\) and setting
\[
  d_R(p,x)=R \qquad (x\in M).
\]
Then \(M_R\) is Lipschitz-free rigid.  In particular, \(M\) embeds isometrically into a Lipschitz-free rigid metric space with one additional point.
\end{theorem}

\begin{proof}
The formula defines a metric: for \(x,y\in M\), one has
\[
  |d_R(p,x)-d_R(p,y)|=0\le d(x,y)
\]
and
\[
  d(x,y)\le \diam(M)<R+R=d_R(x,p)+d_R(p,y).
\]

Because \(M_R\) is finite, a molecule \(m_{a,b}\) is preserved extreme exactly when no third point lies between \(a\) and \(b\).  This follows from the source's extreme-point criterion: in a finite metric space the positive gap in
\[
  d(a,z)+d(z,b)-d(a,b)
\]
over all non-endpoint \(z\) is attained whenever \([a,b]=\{a,b\}\).

We now identify the preserved-extreme graph of \(M_R\).  If \(x,y\in M\), then the new point \(p\) does not lie between them, since
\[
  d_R(x,p)+d_R(p,y)=2R>\diam(M)\ge d(x,y).
\]
Thus \(\{x,y\}\) is preserved extreme in \(M_R\) exactly when it was preserved extreme in \(M\).  On the other hand, if \(x\in M\) and \(y\in M\setminus\{x\}\), then
\[
  d_R(p,x)=R<R+d(x,y)=d_R(p,y)+d(y,x),
\]
and symmetrically \(d_R(p,y)<d_R(p,x)+d(x,y)\).  Hence no old point lies between \(p\) and \(x\), so every pair \(\{p,x\}\) is preserved extreme in \(M_R\).

Consequently \(G_{\mathrm{ext}}(M_R)\) is the cone over \(G_{\mathrm{ext}}(M)\), with cone vertex \(p\).  A cone over a 2-connected graph is 3-connected: after deleting \(p\) and at most one old vertex, the old graph remains connected; if \(p\) is not deleted, then \(p\) is adjacent to every remaining old vertex.  Therefore \(G_{\mathrm{ext}}(M_R)\) is 3-connected.

The space \(M_R\) is finite, and its preserved-extreme graph has all points as vertices, so it is a weak Prague space.  By Cuth--Doucha--Titkos Theorem 4.9 \cite{CDT}, every weak Prague metric space with 3-connected preserved-extreme graph is Lipschitz-free rigid.  Hence \(M_R\) is Lipschitz-free rigid.
\end{proof}

\begin{corollary}
Every finite 2-connected graph, equipped with its graph metric, embeds isometrically into a Lipschitz-free rigid metric space after adding one point.
\end{corollary}

\begin{proof}
For a connected graph with its graph metric, the source Lemma 3.8 gives \(E_{\mathrm{ext}}\) equal to the graph edge set.  Apply Theorem~\ref{thm:one-point}.
\end{proof}

\section*{Verification And Limitations}

The proof gives a one-point positive answer to Question 3 for finite metric spaces whose preserved-extreme graph is 2-connected.  The construction is explicit: put the new point at any common distance \(R>\diam(M)\).

This does not settle the source's full Question 3.  The argument does not cover finite spaces whose preserved-extreme graph has cut vertices, and it does not treat infinite metric spaces.  It also does not address Question 2 on \(\mathbb R^d\).

Novelty check: the run indexes had no exact entry for arXiv:2402.08266.  Web searches for the exact one-point rigidification wording and for ``Lipschitz-free rigid'' variants returned the source paper but no later solution.  Local corpus search found no separate answer.

\begin{thebibliography}{9}
\bibitem{CDT}
M. Cuth, M. Doucha, and T. Titkos,
\emph{Isometries of Lipschitz-free Banach spaces},
arXiv:2402.08266, 2024.
\url{https://arxiv.org/abs/2402.08266}
\end{thebibliography}

\end{document}

